% from ega0-22-fr.tex
\oldpage[0\textsubscript{IV-1}]{196}

\begin{equation}
\label{0.22.4.4.6-fr}
  \xymatrix{
    && 0 \ar[d] & 0 \ar[d] && \\
    & 0 \ar[r] & V_0\otimes_KK' \ar[r] \ar[d] & V_0' \ar[r] \ar[d] & 0 & \\
    0 \ar[r] & N \ar[r] \ar[d] & \Omega^1_A\otimes_AK' \ar[r] \ar[d] &
      \Omega^1_{A'}\otimes_{A'}K' \ar[r] \ar[d] &
      \Omega^1_{A'/A}\otimes_{A'}K' \ar[r] \ar[d] & 0 \\
    0 \ar[r] & \Upsilon_{K'/K} \ar[r] \ar[d] &
      \Omega^1_K\otimes_KK' \ar[r] \ar[d] & \Omega^1_{K'} \ar[r] \ar[d] &
      \Omega^1_{K'/K} \ar[r] \ar[d] & 0 \\
    & 0 & 0 & 0 & 0 &
  }
  \tag{22.4.4.6}
\end{equation}
sont donc formées respectivement des noyaux et des conoyaux des flèches
horizontales médianes, ce qui achève de prouver que le diagramme est
commutatif. En outre~:

\begin{lemma}[22.4.4.7]
\label{0.22.4.4.7-fr}
Les homomorphismes
\[
  N\longrightarrow\Upsilon_{K'/K}\qquad\text{et}\qquad
  \Omega^1_{A'/A}\otimes_{A'}K'\longrightarrow\Omega^1_{K'/K}
\]
du diagramme (22.4.4.6) sont bijectifs.
\end{lemma}

\begin{proof}
En effet, le diagramme du serpent (Bourbaki, \emph{Alg. comm.}, chap.~I,
§~1, n\textsuperscript{o}~4, prop.~2) appliqué aux deux colonnes médianes de
(22.4.4.6) donne une suite exacte
\[
  0\longrightarrow N\longrightarrow\Upsilon_{K'/K}\longrightarrow0
  \longrightarrow\Omega^1_{A'/A}\otimes_{A'}K'
  \longrightarrow\Omega^1_{K'/K}\longrightarrow0.
\]
Or, si $t_i$ est l'image canonique de $T_i$ dans $A'$, la relation
$F_i(t_i)=0$, jointe au fait que dans $A'$ on a
$p\cdot1\in\mathfrak mA'\subset\mathfrak m'$, montre aussitôt que l'on a
$d_{A'}f_i\in\mathfrak m'\Omega^1_{A'}$, donc
$d_{A'}f_i\otimes1_{K'}=0$~; cela prouve que les
$d_Af_i\otimes1_{K'}$ appartiennent à $N$. En outre, leurs images dans
$\Upsilon_{K'/K}$ sont les $d_Ka_i\otimes1_{K'}$, qui forment une base de
$\Upsilon_{K'/K}$ sur $K'$ (21.4.7)~; compte tenu de (22.4.4.7), cela achève
de démontrer (22.4.4).
\end{proof}
\end{proof}
\begin{env}[22.4.5]
\label{0.22.4.5-fr}
Considérons maintenant deux anneaux $A$, $B$ \emph{de caractéristique} $p>0$, et
deux homomorphismes
\[
  A\xrightarrow{i}B\xrightarrow{j}A
\]

